\appendix

\section{A hypergeometric antiderivative}
\label{app:hyp}

The quotient that appears when the \(x\)-characteristic is inserted directly into
the \(y\)-equation has the form
\begin{equation}
  I(\tau)=\int\frac{\e^{h\tau}}
  {a+b\e^{k\tau}}\,\dd\tau.
  \label{eq:Iform}
\end{equation}
An Euler-integral proof would require a sign condition that excludes the
subcritical application.  A local power-series proof gives the correct identity
first, after which analytic continuation supplies the wider domain.

\begin{proposition}
Suppose \(a,h,k\neq0\) and \(h+nk\neq0\) for every integer \(n\geq0\).  On a
simply connected domain where a consistent branch is chosen and
\(a+b\e^{k\tau}\neq0\),
\begin{equation}
  I(\tau)=
  \frac{\e^{h\tau}}{ah}
  {}_2F_1\!\left(
    1,\frac{h}{k};\frac{h}{k}+1;
    -\frac{b}{a}\e^{k\tau}
  \right)+C,
  \label{eq:hypIden}
\end{equation}
where the right-hand side is understood by analytic continuation from its power
series whenever necessary.
\end{proposition}

\begin{proof}
First suppose
\(\lvert(b/a)\e^{k\tau}\rvert<1\).  The geometric series gives
\begin{align}
  \frac{\e^{h\tau}}{a+b\e^{k\tau}}
  &=\frac1a\sum_{n=0}^{\infty}
    \left(-\frac ba\right)^n\e^{(h+nk)\tau},\\
  I(\tau)
  &=\frac1a\sum_{n=0}^{\infty}
    \left(-\frac ba\right)^n
    \frac{\e^{(h+nk)\tau}}{h+nk}+C.
  \label{eq:hyp-series-antiderivative}
\end{align}
Put \(\nu=h/k\).  Since
\[
  {}_2F_1(1,\nu;\nu+1;z)
  =\sum_{n=0}^{\infty}\frac{\nu}{\nu+n}z^n,
\]
the series in \eqref{eq:hyp-series-antiderivative} is precisely
\eqref{eq:hypIden}.  Both sides are analytic away from their poles and chosen
branch cuts, so equality extends throughout any connected domain on which the
same continuation is used.
\end{proof}

If \(h+n_0k=0\) for an integer \(n_0\geq0\), the corresponding term in
\eqref{eq:hyp-series-antiderivative} integrates to
\(
  a^{-1}(-b/a)^{n_0}\tau
\)
rather than an exponential quotient.  This is the secular term that appears at
the resonance \eqref{eq:abd-resonance}; the original finite integral remains
regular.

\section{Recovering state probabilities from a PGF}
\label{app:extract}

A PGF is analytic in the open unit disc, or in the unit polydisc for several
counts.  Its Taylor coefficients are therefore the state probabilities.  For the
bivariate PGF \eqref{pgfDef},
\begin{equation}
  p_{ij}(t)=
  \frac{1}{i!j!}
  \left.
  \frac{\partial^{i+j}G}
  {\partial x^i\partial y^j}
  \right|_{x=y=0}.
  \label{eq:stateProb}
\end{equation}
Analyticity, rather than repeated differentiability alone, is the property that
justifies equality with the full Taylor series.

For the finite-state absorption models,
\begin{equation}
  G(x,y,t)=\bigl[c(t)+f(t)x+g(t)y\bigr]^N.
  \label{eq:multilinear}
\end{equation}
Repeated differentiation gives
\begin{equation}
  \frac{\partial^{i+j}G}{\partial x^i\partial y^j}
  =\frac{N!}{(N-i-j)!}
  f(t)^ig(t)^j
  \bigl[c(t)+f(t)x+g(t)y\bigr]^{N-i-j}
  \label{eq:multilinear-derivative}
\end{equation}
for \(i+j\leq N\), and zero otherwise.  Substitution into
\eqref{eq:stateProb} yields the trinomial law
\begin{equation}
  p_{ij}(t)=
  \frac{N!}{i!j!(N-i-j)!}
  f(t)^ig(t)^jc(t)^{N-i-j},
  \qquad i+j\leq N.
  \label{eq:stateProbSpec}
\end{equation}

For pure absorption, \(c=0\),
\(f=1-\e^{-\alpha t}\) and \(g=\e^{-\alpha t}\), so only states with
\(i+j=N\) have positive probability.  For absorption--death,
\(c=q_{\mathrm d}\), \(f=q_{\mathrm i}\) and \(g=q_{\mathrm o}\), giving the
three independent fates in \eqref{eq:Gabsdeath}.

The absorption--birth--death PGF is not a finite polynomial in \(x\).  Its
coefficients may still be obtained from \eqref{eq:stateProb}; numerically, a
stable alternative is the bivariate Cauchy formula
\begin{equation}
  p_{ij}(t)=\frac{1}{(2\pi\mathrm i)^2}
  \oint_{|x|=r_x}\oint_{|y|=r_y}
  \frac{G(x,y,t)}{x^{i+1}y^{j+1}}\,\dd y\,\dd x,
  \qquad 0<r_x,r_y<1.
  \label{eq:cauchy-coefficient}
\end{equation}
Using the integral representation \eqref{eq:abd-integral-solution} inside these
contours avoids the branch bookkeeping of the hypergeometric form.
